\section*{Chapter \ref{ch:b3:nuclear-physics} --- Nuclear Physics}

\begin{solution}{exo:b3:nuclear-physics:1}
(a) $^{12}$C: $6, 6, 12$; $^{14}$C: $6, 8, 14$; $^{235}$U:
$92, 143, 235$; $^{238}$U: $92, 146, 238$. (b) The carbons; the
uraniums. (c) Chemistry is the electron cloud, fixed by $Z$;
stability is the $N$--$Z$ balance inside. (d) Isobars ---
connected by $\beta^-$ decay: $^{14}\text{C} \to {}^{14}\text{N}
+ \text{e}^- + \bar\nu$, the radiocarbon clock's tick.
\end{solution}

\begin{solution}{exo:b3:nuclear-physics:2}
(a) $\rho = Am_{\text{u}}/\tfrac43\pi r_0^3A =
3m_{\text{u}}/4\pi r_0^3 \approx \qty{2.3e17}{kg/m^3}$: $A$
cancels. (b) $\qty{5}{mL} \to \qty{1.2e12}{kg}$ --- a billion
tonnes in a teaspoon. (c) $(R/R_{\text{atom}})^3 \approx
(4.7\times10^{-5})^3 \approx 10^{-13}$: the atom is empty space
with a heavy speck. (d) Volume $\propto A$ means each nucleon
claims fixed room and binds only its touching neighbours: the
force saturates at \unit{fm} range.
\end{solution}

\begin{solution}{exo:b3:nuclear-physics:3}
(a) $\Delta m = 2(1.007276 + 1.008665) - 4.001506 =
\qty{0.030376}{u}$: $B = \qty{28.3}{MeV}$. (b) $B/A =
\qty{7.07}{MeV}$. (c) $28.3\,\text{MeV}/13.6\,\text{eV} \approx
2\times10^6$: six orders of magnitude --- why nuclear fuel
outclasses chemical by a million. (d) The $\alpha$ is a
pre-assembled, exceptionally cheap parcel: emitting it exports
four nucleons at a bargain \qty{28.3}{MeV} rebate, which no
single nucleon can match.
\end{solution}

\begin{solution}{exo:b3:nuclear-physics:4}
(a) $N = N_{\text{A}}/226 = \qty{2.66e21}{}$. (b) $\lambda =
\ln2/(1600\times\qty{3.156e7}{s}) = \qty{1.37e-11}{s^{-1}}$. (c)
$\mathcal A = \lambda N = \qty{3.7e10}{Bq}$: one curie, by
construction --- the Curies' own gram. (d) Three half-lives:
$\qty{0.125}{g}$ of radium (the rest now radon and its
daughters, marching toward lead).
\end{solution}

\begin{solution}{exo:b3:nuclear-physics:5}
(a) Volume $+882$; surface $-261$; Coulomb $-0.711\times676/3.83
= -126$; asymmetry $-23.7\times16/56 = -6.8$ (\unit{MeV}). (b)
With pairing $+12/\sqrt{56} = +1.6$: $B \approx \qty{490}{MeV}$,
$B/A = 8.76$ against the measured $8.79$ --- three per mille
from a liquid drop. (c) Volume against surface$+$Coulomb: their
crossover in growth rates is what carves the maximum at iron.
(d) Coulomb is long-ranged: every proton repels every other,
$\propto Z^2$, while the saturating strong force only ever pays
$\propto A$ --- the eventual defeat of the heavyweights.
\end{solution}

\begin{solution}{exo:b3:nuclear-physics:6}
(a) $\partial/\partial Z$ of the Coulomb and asymmetry terms
gives $2a_{\text{C}}Z/A^{1/3} = 4a_{\text{A}}(A - 2Z)/A$; solve.
(b) $Z^* = 50.5/1.163 = 43.4$: nature's choice at $A = 101$ is
ruthenium, $Z = 44$. (c) Without Coulomb, $Z^* = A/2$: two Fermi
ladders (\cref{ch:b3:quantum-statistics}), cheapest filled to
equal height. (d) The fragment's valley seat is $Z^* \approx 41$:
five protons short --- a cascade of five $\beta^-$ decays climbs
it back, each converting a neutron and emitting an electron and
antineutrino.
\end{solution}

\begin{solution}{exo:b3:nuclear-physics:7}
(a) $t = T_{1/2}\ln(1/0.3)/\ln2 = 5730\times1.74 \approx
\qty{e4}{}$ years: an end-of-ice-age fire. (b) $N_{^{14}\text{C}}
= (N_{\text{A}}/12)\times10^{-12} = \qty{5.0e10}{}$; $\lambda =
\qty{3.8e-12}{s^{-1}}$: $\mathcal A \approx \qty{0.2}{Bq}$ ---
a dozen decays per minute per gram, the working count rate of
every dating laboratory. (c) After nine half-lives the activity
sinks below background: the clock still runs but can no longer
be read. (d) Rocks never breathed atmospheric carbon; their
clocks are the primordial ones --- uranium--lead,
potassium--argon --- with half-lives of billions of years.
\end{solution}

\begin{solution}{exo:b3:nuclear-physics:8}
(a) The $\alpha$'s \qty{8.78}{MeV} is far under the
\qty{26}{MeV} rim: classically it can neither leave nor have
entered --- yet the half-life is \qty{0.3}{\micro s}. (b)
$\lambda_1 = \ln2/\qty{3e-7}{s} = \qty{2.3e6}{s^{-1}}$;
$\lambda_2 = \ln2/\qty{1.4e17}{s} = \qty{4.9e-18}{s^{-1}}$:
ratio $\approx 5\times10^{23}$. (c) The tunnelling rate is
$\eu^{-2G}$ with the Gamow exponent $\propto Z/\sqrt E$:
halving $E$ raises $2G$ by some fifty-five units, and
$\eu^{55} \approx 10^{24}$ --- Geiger--Nuttall's outrageous
straight line, from one exponential. (d) Solar protons also
tunnel: the core temperature must be just high enough for the
Maxwell tail times the Gamow factor to sustain the burn ---
same cliff, climbed from the other side.
\end{solution}

\begin{solution}{exo:b3:nuclear-physics:9}
(a) The split moves ${\sim}235$ nucleons from \qty{7.6}{} to
${\sim}\qty{8.5}{MeV}$ of binding: $235\times0.9 \approx
\qty{200}{MeV}$. (b) $N = \qty{2.56e24}{}$ nuclei:
$E = \qty{8.2e13}{J} \approx 2000$ tonnes of oil --- per
kilogram. (c) Thermal \qty{3}{GW} for a year is
\qty{9.5e16}{J}: about \qty{1.2}{t} of $^{235}$U. (d) The two
positive fragments spring apart under their own Coulomb
repulsion, carrying ${\sim}\qty{170}{MeV}$ as kinetic energy
that becomes heat within micrometres --- a reactor is a kettle
whose flame is electrostatic recoil.
\end{solution}

\begin{solution}{exo:b3:nuclear-physics:10}
(a) $10^4$ generations: $(1.001)^{10^4} = \eu^{10} \approx
22000$ --- megawatts to tens of gigawatts inside a second; no
motor moves a rod that fast. (b) With $k - 1 < \beta$ the chain
cannot reproduce on prompt neutrons alone: every growth step
waits for the ${\sim}\qty{10}{s}$ stragglers, so the effective
generation time is a fraction of a second or more. (c)
$(1.0005)^{10} \approx 1.005$: half a percent per second ---
rod-and-operator territory. (d) Driving $k - 1$ past $\beta$
put the reactor on the \qty{e-4}{s} prompt clock: they lost the
delayed-neutron grace period that makes reactors steerable.
\end{solution}

\begin{solution}{exo:b3:nuclear-physics:11}
(a) $\dot m = L_\odot/c^2 = \qty{4.2e9}{kg/s}$ of pure mass;
hydrogen throughput $\dot m/0.0071 \approx \qty{6e11}{kg/s}$ ---
six hundred million tonnes a second. (b) Burnable hydrogen
${\sim}2\times10^{30}\times0.1\times0.75 = \qty{1.5e29}{kg}$:
$t \approx \qty{2.5e17}{s} \approx 8$ billion years --- the Sun
is middle-aged. (c) Barrier ${\sim}\qty{1.4}{MeV}$ against
$k_{\text{B}}T \approx \qty{1.3}{keV}$: a thousandfold deficit
--- only the Maxwell tail's fastest protons, tunnelling
(\cref{exo:b3:nuclear-physics:8}), ever fuse; hence the Sun
burns for gigayears instead of exploding. (d) The ash is
$^4$He --- stable, doubly magic contentment; there are no
neutron-rich fragments to $\beta$-decay for centuries.
\end{solution}

\begin{solution}{exo:b3:nuclear-physics:12}
(a) $\lambda = \ln2/\qty{21600}{s} = \qty{3.2e-5}{s^{-1}}$:
$N = \mathcal A/\lambda = \qty{1.6e13}{}$ nuclei ---
\qty{2.6}{ng}. Medicine by the nanogram. (b) Six hours outlives
the scan but not the weekend: the dose switches itself off; a
pure \qty{140}{keV} $\gamma$ leaves the body to the camera
without the tissue-burning $\alpha$/$\beta$ toll. (c) Electron
and positron at rest annihilate into \emph{two}
\qty{511}{keV} photons, back to back (energy and momentum,
\cref{ch:b3:relativistic-dynamics}): each coincident pair
defines a line through the tumour, and thousands of lines
triangulate it --- tomography by conservation law. (d)
\qty{2}{J/kg} warms tissue half a millikelvin --- thermally a
sip of tepid water; but delivered as ${\sim}10^{17}$ ionisations
per kilogram, each snipping \unit{eV}-scale chemical bonds, it
is chemistry-lethal to a dividing cell. Joules measure heat;
grays measure sabotage.
\end{solution}

\begin{solution}{pb:b3:nuclear-physics:1}
\textbf{1.} $\qty{200}{MeV} = \qty{3.2e-11}{J}$:
$20\times10^6/3.2\times10^{-11} = \qty{6.3e17}{}$ fissions per
second.
\textbf{2.} $5.4\times10^{22}$ fissions per day $\times
235/N_{\text{A}}$: about \qty{21}{g} of $^{235}$U a day.
\textbf{3.} Coal: $\qty{1.7e12}{J}/\qty{3e7}{J/kg} \approx 58$
tonnes a day --- a mass ratio of $\sim3\times10^{6}$, which
\emph{is} the \unit{MeV}-to-\unit{eV} ratio of nuclear to
chemical bonds.
\textbf{4.} $235\times(8.5 - 7.6) \approx \qty{210}{MeV}$:
consistent, the small change bookkept by neutrons and
neutrinos.
\textbf{5.} Fragment kinetic energy $\to$ ionisation $\to$ heat,
within micrometres and microseconds; conduction hands it to the
water --- the reactor is a fragment-stopping kettle.
\textbf{6.} They inherit uranium's $N/Z \approx 1.55$, far
above the valley floor at their $A$: each must run a chain of
$\beta^-$ decays back down --- built-in radioactivity, by
geometry of the valley.
\textbf{7.} $k$ = neutrons begotten per neutron, one generation
on: $0.999$ dies out, $1.000$ holds steady (an operating
reactor), $1.001$ grows.
\textbf{8.} Slow neutrons linger near the nucleus and
$^{235}$U's fission appetite grows steeply at thermal energies.
Elastic collisions shed energy fastest onto equal masses ---
and hydrogen matches the neutron's mass: water is moderator
made to order.
\textbf{9.} Boron-10 devours neutrons ($\text{n} + {}^{10}
\text{B} \to \alpha + {}^{7}\text{Li}$): rods in, neutrons
eaten, $k$ down; rods out, $k$ up --- the throttle.
\textbf{10.} $(1.0005)^{10^4} = \eu^{5} \approx 150$: powers of
a hundred and fifty per second --- hopeless for machinery.
\textbf{11.} On the \qty{0.1}{s} delayed clock:
$(1.0005)^{10} \approx 1.005$, half a percent per second. Rule:
keep $|k - 1|$ well below $\beta = 0.007$, so the delayed
neutrons always hold the casting vote.
\textbf{12.} Boiling removes the moderator: neutrons stay fast,
fission starves, $k$ falls --- the water-moderated design
throttles itself (a negative feedback its graphite cousins
lacked).
\textbf{13.} Decay heat: the fragments' radioactivity ---
Part I question 6 --- obeys half-lives, not control rods, and
still yields megawatts just after shutdown.
\textbf{14.} Outrun light in the water: $v > c/n$.
\textbf{15.} $v > 0.752c$: $\gamma = 1.51$, so $E_{\text{k}} >
0.51\times\qty{511}{keV} \approx \qty{0.26}{MeV}$.
\textbf{16.} $\beta$ electrons from fragment decays carry
\unit{MeV}s, and core $\gamma$s Compton-kick electrons to
similar energies: both sail past \qty{0.26}{MeV} --- the pool
is full of qualifying electrons.
\textbf{17.} The Cherenkov spectrum strengthens toward short
wavelengths ($\propto1/\lambda^2$ in intensity): the eye is
handed blue and violet --- the colour is the spectrum's slope.
\textbf{18.} Water attenuates $\gamma$s and neutrons
exponentially; eight metres is dozens of halving-lengths, so
the balcony sits at background dose. The glow is the water
\emph{absorbing} the radiation's energy --- visible proof the
shield is eating it.
\textbf{19.} The short-lived fragments die within minutes
(the dimming); the longer-lived tail --- the same decay heat of
question 13 --- keeps a faint glow and a real heat load for
days.
\textbf{20.} With $T_{1/2} = \qty{66}{h}$ the parent loses a
factor ${\sim}6$ per week: stockpiles decay themselves away,
so hospitals live on a weekly ``technetium cow'' delivered from
reactors like this one.
\textbf{21.} $96/66 = 1.45$ half-lives: $2^{-1.45} \approx
0.36$ --- about \qty{36}{GBq} remain.
\textbf{22.} Thermalised to room temperature, $\lambda =
h/\sqrt{3m_{\text{n}}k_{\text{B}}T} \approx \qty{1.5}{\angstrom}$
(\cref{exo:b3:crystalline-solids:5}): born matched to \omterm{def:b3:crystalline-solids:lattice}{lattice}
spacings.
\textbf{23.} Five years kills every half-life up to months:
activity and decay heat fall by orders of magnitude, until the
element can ride in a shielded cask instead of a swimming pool.
\textbf{24.} ${\sim}\qty{200}{kW}$ --- a hundred domestic
heaters into a hundred-tonne pool: tens of kelvin per day at
worst, comfortably removed by natural convection --- passive
safety by sheer heat capacity.
\textbf{25.} $B/A$ pays \qty{200}{MeV} a split and \qty{21}{g}
a day; delayed neutrons lend the seconds that make $k$
steerable; the water moderates, cools, shields --- and glows
blue precisely because it is working; and the product that
matters most leaves in vials for Monday's hospitals.
\end{solution}
